\documentclass[10pt,a4paper]{amsart}
\usepackage[margin=19mm]{geometry}
\usepackage{amsmath,amssymb,mathtools}
\usepackage[scr=rsfs,cal=euler]{mathalfa}
\usepackage{xfrac}
\usepackage[unicode,hidelinks,bookmarks=false]{hyperref}
\newcommand{\ie}{i.e.\ }
\newcommand{\cf}{cf.\ }
\newcommand{\resp}{respectively\ }
\newcommand{\Subsection}{\subsection}
\providecommand{\nobreakdash}{\nobreak\hbox{-}}
\setlength{\emergencystretch}{2em}
\multlinegap0pt
\usepackage{color}
\usepackage{epstopdf}
\usepackage{algorithmicx}
\usepackage{booktabs}
\usepackage{latexsym}
\usepackage{mathtools}
\usepackage{algorithm}
\usepackage{algpseudocode}
\usepackage{epstopdf}
\usepackage{upref}
\usepackage{fancyvrb}
\usepackage{mathtools}
\usepackage{float}
\usepackage{subfig}
\usepackage{commath}
\usepackage{subfig}
\usepackage{placeins}
\newenvironment{ecc}{\begin{quote}\small\sf{\color{green}$\diamondsuit$~ECC~}}{\end{quote}}
\newenvironment{dk}{\begin{quote}\small\sf{\color{red}$\clubsuit$~DK~}}{\end{quote}}
\newcommand{\What}{\widehat{W}}
\newcommand{\Shat}{\widehat{S}}
\newcommand{\Vhat}{\widehat{V}}
\newcommand{\Mhat}{\widehat{M}}
\newcommand{\Lhat}{\widehat{L}}
\newcommand{\Phat}{\widehat{P}}
\newcommand{\Qhat}{\widehat{Q}}
\newcommand{\That}{\widehat{T}}
\newcommand{\LhatT}{\widehat{L}^T}
\newcommand{\gammap}{\gamma_p}
\newcommand{\gammaxi}{\gamma_\xi}
\newcommand{\E}{\theta}
\newcommand{\F}{\phi}
\def\xik#1{\xi^{(#1)}}
\def\lambdak#1{\lambda^{(#1)}}
\def\x#1{x^{(#1)}}
\def\xT#1{x^{(#1)T}}
\def\xhat#1{\widehat{x}^{(#1)}}
\def\xhatT#1{\widehat{x}^{(#1)T}}
\def\r#1{r^{(#1)}}
\def\rhat#1{\widehat{r}^{(#1)}}
\def\d#1{d^{(#1)}}
\def\dhat#1{\widehat{d}^{(#1)}}
\def\normo#1{\|#1\|_{\infty}}
\def\normF#1{\|#1\|_F}
\def\norm#1{\|#1\|}
\def\normt#1{\|#1\|_2}
\def\normn#1{\|#1\|_1}
\def\normA#1{\|#1\|_A}
\def\normAinv#1{\|#1\|_{A^{-1}}}
\newtheorem{theorem}{Theorem}[section]
\newtheorem{lemma}[theorem]{Lemma}
\newtheorem{remark}[theorem]{Remark}
\newtheorem{corollary}{Corollary}[theorem]
\begin{document}
\title[Error Analysis and Precision Selection for Mixed-Precision D]{Error Analysis and Precision Selection for Mixed-Precision DEIM-CUR Decompositions}
\author{Ioannis Thanasis, Erin Carson}
\date{}
\maketitle
\noindent\emph{Source and licence.} Error Analysis and Precision Selection for Mixed-Precision DEIM-CUR Decompositions. Version arXiv:2609.24509v2 (CC BY 4.0). This material is distributed under the Creative Commons Attribution 4.0 International licence, http://creativecommons.org/licenses/by/4.0/, as recorded in the arXiv OAI licence metadata for this exact version and shown on the abstract page https://arxiv.org/abs/2609.24509v2. Full rights notice: licenses/arxiv-2609.24509-CC-BY-4.0.txt.\par\smallskip
\noindent\textit{Editorial scope.} Counted unit: the original \texttt{theorem} environment \texttt{thm:main} at source lines 280--294 together with its complete proof at lines 295--297; the delivered document keeps source lines 177--298 so that every referenced label and every citation resolves inside it. Native editable source is the paper's own concatenated TeX; the AMS wrapper is editorial formatting only. No publisher class, upstream script or unrelated artwork is redistributed.
\par\smallskip\noindent\textit{Reference note.} This article ships no compiled bibliography (biblatex); the entries delivered here are composed from the author's own BibTeX (.bib) fields, with no field invented and no entry dropped.
\section{CUR Approximation with Inexact SVD}
\label{sec:inexactSVD}

\hspace{5mm}Let $A = V S W^T$ be the exact full SVD of $A$, and let $A_k = V_k S_k W_k^T$ be the best rank-$k$ approximation of $A$ given by the (exactly computed) truncated SVD of $A$, where $V_k \in \mathbb{R}^{m\times k}, S_k \in \mathbb{R}^{k \times k}$, and $W_k \in \mathbb{R}^{n \times k}$. 

Assume that we have computed a rank-$k$ truncated SVD of an $m\times n$ matrix $A$ inexactly in some way, and that the inexact SVD is given by $\widehat{A}_k = \Vhat_k \Shat_k \What_k^T$. Thus $\Shat_k$ does not contain the exact leading $k$ singular values, and we do not assume that $\Vhat_k$ and $\What_k$ have orthonormal columns. 
We assume that we can write the bound
 \begin{equation}
 \norm{A-\Vhat_k \Shat_k \What_k^T} \leq \E \sigma_{k+1} +\F,
 \label{eq:inexactSVD}
 \end{equation}
 % 
for some quantities $\E$ and $\F$. The reason for this format will become clear when we later apply our results to specific inexact SVD computations. Note that for the exact SVD computation, we have $\E=1$ and $\F=0$.


 We consider a CUR factorization that uses row indices $p\in\mathbb{N}^k$ and column indices $q\in\mathbb{N}^k$, and set
\begin{equation}
P=I(:,p)\in\mathbb{R}^{m\times k}, \quad Q=I(:,q)\in\mathbb{R}^{n\times k}.
\label{eq:PQdef}
\end{equation}
We first prove an analog of \cite[Lemma 4.1]{sorensen2016deim} without the assumption that $\Vhat_k^T \Vhat_k = I$. 

\begin{lemma}
Let $\widehat{A}_k = \Vhat_k \Shat_k \What_k^T$ be an inexactly computed rank-$k$ truncated SVD of $A$ satisfying the bound \eqref{eq:inexactSVD}, and let $P$ and $Q$ be defined as in \eqref{eq:PQdef}.
Assume that $P^T\Vhat_k$ is invertible, and let $\mathcal{P}=\Vhat_k(P^T\Vhat_k)^{-1}P^T$ be an interpolatory projector. Then 
%
\begin{equation*}
\norm{A-\mathcal{P}A} \leq \norm{\Vhat_k} \norm{(P^T\Vhat_k)^{-1}} (\E\sigma_{k+1}+\F).
\end{equation*}

\end{lemma}
\begin{proof}
    
 Let $\mathcal{P}=\Vhat_k(P^T\Vhat_k)^{-1}P^T$ be an interpolatory projector. From this we have
\[
\mathcal{P}\Vhat_k = \Vhat_k \rightarrow (I-\mathcal{P})\Vhat_k = 0.
\]
From this,
\begin{align*}
\norm{A-\mathcal{P}A} = \norm{(I-\mathcal{P})A} &= \norm{(I-\mathcal{P})(I-\Vhat_k \Vhat_k^\dagger)A} \\
&\leq \norm{I-\mathcal{P}} \norm{(I-\Vhat_k \Vhat_k^\dagger)A}.
\end{align*}
We can then write, assuming $\mathcal{P}\neq 0$ or $I$,
\[
\norm{I-\mathcal{P}} = \norm{\mathcal{P}} = \norm{\Vhat_k (P^T\Vhat_k)^{-1}P^T} = \norm{\Vhat_k (P^T\Vhat_k)^{-1}} \leq \norm{\Vhat_k}\norm{(P^T \Vhat_k)^{-1}};
\]
note that this differs from the result in \cite[Lemma 4.1]{sorensen2016deim} since we cannot assume $\norm{\Vhat_k}=1$.
Thus altogether we have the bound
\[
\norm{A-\mathcal{P}A} \leq \norm{\Vhat_k} \norm{(P^T\Vhat_k)^{-1}} \norm{(I-\Vhat_k\Vhat_k^\dagger)A}.
\]

We now seek to bound the last term on the right-hand side above. Writing $A = \widehat{A}_k + (A-\widehat{A}_k)$, we have
\begin{equation}
    \norm{(I-\Vhat_k\Vhat_k^\dagger)A} = \norm{(I-\Vhat_k \Vhat_k^\dagger)(A-\widehat{A}_k)} \leq \norm{A-\widehat{A}_k} \leq \E\sigma_{k+1} + \F.
\end{equation}



Substituting this above, we have 
\[
\norm{A-\mathcal{P}A} \leq \norm{\Vhat_k} \norm{(P^T\Vhat_k)^{-1}} (\E\sigma_{k+1}+\F).
\]

\end{proof}

In a completely analogous way, we can consider the right interpolatory projector $\mathcal{Q} = Q(\What_k^TQ)^{-1}\What_k^T$, assuming that $\What^TQ$ is invertible, and can show that 
\begin{equation}
    \Vert A - A\mathcal{Q}\Vert \leq \norm{\What_k}\norm{(\What_k^T Q)^{-1}}(\E\sigma_{k+1}+\F).
\end{equation}


Now, define
\begin{equation}
    \tilde{\eta}_p = \norm{(P^T\Vhat_k)^{-1}}, \qquad 
    \tilde{\eta}_q = \norm{(\What_k^T Q)^{-1}}.
    \label{eq:etas}
\end{equation}
Then 
\begin{align}
    \norm{(I-\mathcal{P})A} & \leq \norm{\Vhat_k}\tilde{\eta}_p(\E\sigma_{k+1}+\F),\\
    \norm{A(I-\mathcal{Q})} &\leq \norm{\What_k}\tilde{\eta}_q (\E\sigma_{k+1}+\F).
\end{align}

The CUR approximation uses $C = AQ$ and $R = P^T A$.
The relevant projection residuals involve the orthogonal projectors
$CC^\dagger$ and $R^\dagger R$.
\begin{lemma}
\label{lem:CR}
Assuming \eqref{eq:inexactSVD} holds, for $C = AQ$ and $R = P^T A$ with $Q$ and $P$ defined in \eqref{eq:PQdef},
\begin{align}
  \norm{(I - CC^\dagger)A} &\leq \norm{\What_k}\,\tilde\eta_q\,(\E\sigma_{k+1}+\F),
  \label{eq:C_res} \\
  \norm{A(I - R^\dagger R)} &\leq \norm{\Vhat_k}\,\tilde\eta_p\,(\E\sigma_{k+1}+\F).
  \label{eq:R_res}
\end{align}
\end{lemma}

\begin{proof}
The proof is analogous to the proof of Lemma 4.2 in \cite{sorensen2016deim}.
\end{proof}

We can then state the result for the error of the CUR decomposition. 

\begin{theorem}
\label{thm:main}
Let $A \in \mathbb{R}^{m\times n}$, $1 \le k < \min(m,n)$.  Suppose the DEIM index
sets $P$ and $Q$ are computed from $\Vhat_k$ and $\What_k$
satisfying \eqref{eq:inexactSVD}, with amplification factors
$\tilde\eta_p$ and $\tilde\eta_q$ defined in \eqref{eq:etas}.  Let $C = AQ$, $R = P^T A$, and
$U = C^\dagger A R^\dagger$.  Then
\begin{equation}
  \norm{A - CUR}
  \leq
  \bigl(\norm{\What_k}\,\tilde\eta_q + \norm{\Vhat_k}\,\tilde\eta_p\bigr)
  \bigl(\E\sigma_{k+1} + \F\bigr).
  \label{eq:main}
\end{equation} 
\end{theorem}

\begin{proof}
    The proof is analogous to the proof of Theorem 4.1 in \cite{sorensen2016deim}.
\end{proof}

\begin{thebibliography}{99}
\bibitem[]{sorensen2016deim} Sorensen, Danny C; Embree, Mark. A {DEIM} induced {CUR} factorization. SIAM Journal on Scientific Computing 38 (2016) A1454--A1482
\end{thebibliography}
\end{document}
